% corpus_supplement: Erdős problem #375
% source: https://www.erdosproblems.com/latex/375
% fetched_at: 2026-07-14T18:48:14Z
% payload_sha256: 2d8583414cc098b31002f4451b326a85d462e1df0c484a781149c841e2e8a979
% statement/remarks/references extracted by erdos_ingest.extract_source_page;
% rendered by erdos_ingest.render_tex — do not edit
\documentclass{article}
\usepackage{amsmath, amssymb, amsthm}
\usepackage{hyperref}
\usepackage{geometry}
\geometry{margin=1in}

\title{Erd\H{o}s Problem \#375}
\date{Last updated: 2025-08-31}
\author{Source: erdosproblems.com}

\begin{document}
\maketitle

\noindent\textbf{Status:} OPEN \quad \textbf{No prize}

\noindent\textbf{Tags:} number theory

\medskip
\noindent\textbf{Problem Statement:}

Is it true that for any $n,k\geq 1$, if $n+1,\ldots,n+k$ are all composite then there are distinct primes $p_1,\ldots,p_k$ such that $p_i\mid n+i$ for $1\leq i\leq k$?

\bigskip
\noindent\textbf{Remarks:}

Note this is trivial when $k\leq 2$. Originally conjectured by Grimm \cite{Gr69}. This is a very difficult problem, since it in particular implies $p_{n+1}-p_n <p_n^{1/2-c}$ for some constant $c>0$, in particular resolving Legendre's conjecture.

Grimm proved that this is true if $k\ll \log n/\log\log n$. Erd\H{o}s and Selfridge improved this to $k\leq (1+o(1))\log n$. Ramachandra, Shorey, and Tijdeman \cite{RST75} have improved this to\[k\ll\left(\frac{\log n}{\log\log n}\right)^3.\]Laishram and Shorey \cite{LaSh06} have verified this is true, for all $k\geq 1$, for all $n\leq 1.9\times 10^{10}$. 

This is problem B32 in Guy's collection \cite{Gu04}.

See also [860].

\bigskip
\noindent\textbf{References:}

[Gr69] Grimm, C. A., A conjecture on consecutive composite numbers. Amer. Math. Monthly (1969), 1126--1128.
 
 [Gu04] Guy, Richard K., Unsolved problems in number theory. (2004), xviii+437.
 
 [LaSh06] Laishram, Shanta and Shorey, T. N., Grimm's conjecture on consecutive integers. Int. J. Number Theory (2006), 207--211.
 
 [RST75] Ramachandra, K. and Shorey, T. N. and Tijdeman, R., On Grimm's problem relating to factorisation of a block of consecutive integers. J. Reine Angew. Math. (1975), 109-124.

\bigskip
\noindent\small{Source: \url{https://www.erdosproblems.com/375}}
\end{document}
